\chapter{Protections and Mitigations}
\label{ch:protections}

\epigraph{``Every mitigation raises the cost of exploitation. None of them makes exploitation impossible.''}{}

\learninggoal{
	\item Describe the major defensive mechanisms in modern operating systems and compilers.
	\item Explain how each mitigation disrupts specific exploitation techniques.
	\item Understand the interaction between mitigations and the exploitation pipeline.
	\item Compare hardware, compiler, and OS-level protection strategies.
	\item Evaluate the effectiveness and limitations of current mitigations.
}

\section{The Mitigation Stack}

Modern systems deploy defences at multiple layers:

\begin{longtable}{p{3cm}p{4cm}p{5cm}}
	\toprule
	\textbf{Layer} & \textbf{Mitigation} & \textbf{What it prevents}              \\
	\midrule
	Hardware       & NX bit              & Direct shellcode execution             \\
	Hardware       & SMEP/SMAP           & Kernel-mode code reuse from user pages \\
	OS             & ASLR                & Predictable address computation        \\
	OS             & KASLR               & Kernel address prediction              \\
	Compiler       & Stack canary        & Stack buffer overflow detection        \\
	Compiler       & CFI                 & Arbitrary control flow hijacking       \\
	Compiler       & SafeStack           & Separate stack for sensitive data      \\
	Runtime        & AddressSanitizer    & Heap/stack/bounds violation detection  \\
	Runtime        & seccomp             & System call filtering                  \\
	\bottomrule
\end{longtable}

\section{Non-Executable Memory (NX/DEP)}

The NX (No-Execute) bit, called Data Execution Prevention (DEP) on Windows, marks memory pages as non-executable. This prevents the classic exploitation technique of writing shellcode to the stack or heap and executing it.

\begin{definition}[NX/DEP]
	\textbf{NX} is a hardware feature that allows each page table entry to specify whether the page is executable. The processor checks this bit on every instruction fetch; if \texttt{NX=1} and the processor is not in supervisor mode, a page fault is raised.
\end{definition}

NX is nearly universal on modern systems. It is not a complete defence---attackers use code-reuse techniques (ret2libc, ROP) instead---but it eliminates the simplest exploitation path.

\begin{remark}
	Some JIT compilers (e.g., JavaScript JITs in browsers) need executable memory at runtime. They use \texttt{mprotect} or \texttt{VirtualProtect} to mark pages as executable, creating a window of opportunity if an attacker can corrupt JIT code.
\end{remark}

\section{Address Space Layout Randomisation (ASLR)}

\ASLR{} randomises the base address of each memory region (stack, heap, libraries, executable) when a process starts. This prevents the attacker from knowing the absolute address of gadgets, shellcode, or target buffers.

\begin{definition}[ASLR]
	\textbf{Address Space Layout Randomisation} is a technique that randomises the virtual address at which each memory region is mapped. The randomness is typically on the order of bits (e.g., 28 bits for the mmap base on x86-64 Linux), making it infeasible to guess the correct address.
\end{definition}

\subsection{ASLR Effectiveness}

ASLR is effective against attacks that require absolute addresses. The probability of a successful guess is:

\[
	P_{\text{success}} = \frac{1}{2^{N}}
\]

where \(N\) is the number of entropy bits. On x86-64 Linux:

\begin{itemize}
	\item Stack: 22 bits of entropy (aligned to 16 bytes).
	\item mmap (shared libraries): 28 bits.
	\item Heap: 13 bits.
	\item Executable: 0 bits by default (but PIE executables add 28 bits).
\end{itemize}

A brute-force attempt that successfully guesses an address with $2^{28}$ possibilities on the first try has probability $3.7\times 10^{-9}$.

\subsection{Bypassing ASLR}

Attackers bypass ASLR through \emph{information leaks}: any vulnerability that reveals a pointer or code address. Common leaks:

\begin{itemize}
	\item Format string bugs (\texttt{\%p}) leak stack and code addresses.
	\item Heap overflows that corrupt metadata and leak pointers through allocator operations.
	\item Side-channel attacks (Spectre) that leak address information through cache timing.
	\item Uninitialised memory that contains stale pointer values.
\end{itemize}

\begin{principle}[ASLR + Info Leak = No ASLR]
	ASLR is only as strong as the information hiding that supports it. A single leaked pointer renders ASLR ineffective for the entire address-space region containing that pointer.
\end{principle}

\subsection{Position-Independent Executables (PIE)}

A Position-Independent Executable (PIE) compiles the main binary as position-independent code, enabling ASLR for the executable itself. Without PIE, the executable's code is at a fixed address, providing stable targets for ROP gadgets.

\begin{lstlisting}[caption=Checking if a binary is PIE]
  $ readelf -h /bin/ls | grep Type
  Type: DYN (Shared object file)    # PIE
  $ readelf -h /bin/cat | grep Type
  Type: EXEC (Executable file)       # Not PIE (legacy)
\end{lstlisting}

Most modern distributions compile all executables as PIE.

\section{Stack Canaries}

A stack canary is a random value placed between local variables and the saved return address on the stack. Before the function returns, the canary is checked. If it has changed, a buffer overflow has occurred and the program aborts.

\begin{definition}[Stack Canary]
	A \textbf{stack canary} is a secret value written into the stack frame during the function prologue and verified before the epilogue. Corruption of the canary indicates a stack buffer overflow.
\end{definition}

\begin{lstlisting}[language=C, caption=Canary-protected stack frame]
void vulnerable(char *input) {
    char buf[64];
    // Compiler inserts:
    // canary = __readgsqword(0x28);  // read from TLS
    // ... store canary on stack ...
    strcpy(buf, input);
    // ... check canary before ret ...
    // if (canary != __stack_chk_guard) __stack_chk_fail();
}
\end{lstlisting}

The canary is:
\begin{itemize}
	\item Random: different per process, different per thread.
	\item Secret: stored in thread-local storage (TLS).
	\item Terminated with a null byte: prevents string-based leaks.
\end{itemize}

The null-byte terminator means that a canary byte cannot be \texttt{0x00}. If an overflow is done byte-by-byte (e.g., one byte past a buffer), the attacker would overwrite the null byte, which is detected. More importantly, the null byte prevents \texttt{strcpy} and similar functions from reading past the canary.

\subsection{Bypassing Stack Canaries}

Stack canaries are bypassed by:

\begin{itemize}
	\item \textbf{Information leaks:} if the canary value is leaked (e.g., through a format string bug), the attacker can include the correct canary in the overflow payload.
	\item \textbf{Overwrite-before-check:} some code paths allow overwriting a function pointer \emph{before} the return (e.g., overflow within the same frame that corrupts a local function pointer).
	\item \textbf{Non-stack overflows:} canaries only protect stack frames. Heap overflows, use-after-free, and format string bugs are unaffected.
\end{itemize}

\section{Control-Flow Integrity (CFI)}

Control-Flow Integrity~\cite{abadi2005cfi} restricts the set of targets to which indirect branches (function pointers, vtables, return addresses) can point.

\begin{definition}[CFI]
	\textbf{Control-Flow Integrity} ensures that every indirect control transfer (call, jump, return) goes to a target that was valid at compile time. A \emph{forward-edge} CFI constrains calls and jumps; a \emph{backward-edge} CFI constrains returns.
\end{definition}

\subsection{Forward-Edge CFI}

Forward-edge CFI checks that indirect calls and jumps target a valid set of addresses. The compiler determines, for each call site, the set of functions that could legally be reached (based on type information and control-flow analysis). At runtime, before the indirect branch, the target is validated against this set.

Implementation approaches:

\begin{itemize}
	\item \textbf{Clang CFI:} uses type-based analysis to restrict indirect call targets at link time. Overhead is typically 1--5\%.
	\item \textbf{Microsoft Control Flow Guard (CFG):} checks indirect call targets against a bitmap of valid function addresses. Overhead is minimal.
	\item \textbf{Intel CET (Control-flow Enforcement Technology):} hardware-assisted CFI using a shadow stack (backward-edge) and indirect branch tracking (IBT, forward-edge).
\end{itemize}

\subsection{Backward-Edge CFI (Shadow Stacks)}

A shadow stack is a separate, protected memory region that stores a copy of each return address. On function return, the processor compares the return address on the data stack with the shadow stack:

\begin{lstlisting}[caption=Shadow stack concept]
  ; Function prologue
  push return_addr      ; on normal stack
  push return_addr      ; on shadow stack (protected)

  ; Function epilogue
  pop shadow_addr       ; from shadow stack
  pop return_addr       ; from normal stack
  if (shadow_addr != return_addr) fault();
\end{lstlisting}

Shadow stacks prevent ROP (which rewrites return addresses on the data stack) but require hardware support (Intel CET Shadow Stack) or significant software overhead (software shadow stacks like \texttt{retpoline}).

\section{Compiler Hardening}

Modern compilers offer numerous security hardening options:

\begin{longtable}{p{4cm}p{4cm}p{4cm}}
	\toprule
	\textbf{Flag}                     & \textbf{Effect}                                    & \textbf{Overhead} \\
	\midrule
	\texttt{-fstack-protector}        & Canary for functions with local buffers            & Minimal           \\
	\texttt{-fstack-protector-strong} & Canary for any function with local array           & Slightly more     \\
	\texttt{-fstack-clash-protection} & Prevent stack clash probing                        & Small             \\
	\texttt{-D\_FORTIFY\_SOURCE=2}    & Bounds check for \texttt{\_\_builtin\_memcpy} etc. & Near zero         \\
	\texttt{-fPIE -pie}               & Position-independent executable                    & Minimal           \\
	\texttt{-Wl,-z,relro}             & Read-only relocations after loading                & None              \\
	\texttt{-Wl,-z,now}               & Full RELRO (resolve at load time)                  & Startup cost      \\
	\texttt{-fsanitize=address}       & Full runtime bounds + UAF detection                & 2--3x slowdown    \\
	\bottomrule
\end{longtable}

\section{Runtime Defences}

\subsection{seccomp}

seccomp (secure computing mode) is a Linux kernel feature that restricts which system calls a process can make. A process can install a seccomp filter (written as a BPF program) that is evaluated on every \texttt{syscall}:

\begin{lstlisting}[language=C, caption=Minimal seccomp filter]
#include <seccomp.h>

scmp_filter_ctx ctx = seccomp_init(SCMP_ACT_KILL);
seccomp_rule_add(ctx, SCMP_ACT_ALLOW, SCMP_SYS(read), 0);
seccomp_rule_add(ctx, SCMP_ACT_ALLOW, SCMP_SYS(write), 0);
seccomp_rule_add(ctx, SCMP_ACT_ALLOW, SCMP_SYS(exit), 0);
seccomp_load(ctx);
\end{lstlisting}

Even if an attacker achieves arbitrary code execution, seccomp restricts what they can do. For example, blocking \texttt{execve} prevents shell spawning; blocking \texttt{mprotect} prevents making NX pages executable.

\subsection{Sanitizers}

Sanitizers are runtime instrumentation tools that detect memory safety violations:

\begin{itemize}
	\item \textbf{AddressSanitizer (ASan):} detects heap and stack buffer overflows, use-after-free, and double free. Uses shadow memory to track valid memory regions.
	\item \textbf{MemorySanitizer (MSan):} detects reads of uninitialised memory.
	\item \textbf{UndefinedBehaviorSanitizer (UBSan):} detects various forms of undefined behaviour (integer overflow, misaligned access, etc.).
	\item \textbf{LeakSanitizer (LSan):} detects memory leaks.
\end{itemize}

Sanitizers impose runtime overhead (2--5x for ASan) and are intended for testing, not production deployment.

\section{The Arms Race}

Mitigations and bypass techniques co-evolve:

\begin{longtable}{p{3.5cm}p{4.5cm}p{4.5cm}}
	\toprule
	\textbf{Mitigation} & \textbf{Bypass}           & \textbf{Countermeasure}                       \\
	\midrule
	NX/DEP              & ROP                       & CFI, shadow stack                             \\
	ASLR                & Information leak          & Efficient randomization, zero-leak primitives \\
	Stack canary        & Thread-local storage leak & SafeStack, canary rotation                    \\
	PIE                 & Partial overwrite         & ASLR with high entropy                        \\
	CFI (coarse)        & Out-of-bounds target call & Fine-grained CFI, CET IBT                     \\
	seccomp             & Syscall number mutation   & BPF verification of syscall args              \\
	KASLR               & Prefetch side channel     & KPTI, FG-KASLR                                \\
	\bottomrule
\end{longtable}

\begin{principle}[Defence in Depth]
	No single mitigation is sufficient. The goal of defence in depth is to require multiple independent bypasses for a successful exploit. A ROP chain bypasses NX but requires an ASLR leak; a CFI bypass requires a different vulnerability class.
\end{principle}

\section{Exercises}

\begin{exercise}
	Given a binary compiled with \texttt{-fstack-protector} but no PIE and no ASLR, design an exploitation strategy. Which bypass techniques are unnecessary? Which remain?
\end{exercise}

\begin{exercise}
	Ubuntu 22.04 ships with \texttt{-fstack-protector-strong} and PIE by default. A local service has a heap overflow in a function that does not use stack buffers. Which mitigations stop the attacker? Which are irrelevant?
\end{exercise}

\begin{exercise}
	Read about Intel CET. What two components does it provide? How does the shadow stack prevent ROP, and why does it require a new processor instruction?
\end{exercise}

\begin{exercise}
	Design a seccomp filter for a web server. Which syscalls must be allowed? Which can be safely blocked? Would your filter prevent an attacker who achieves RCE from spawning a reverse shell?
\end{exercise}
